\section{Công trình liên quan}
\label{sec:related-work}

\TBD{}

% ==== BẢN NHÁP CŨ (giữ tham khảo, sẽ viết lại) ====
% \section{Công trình liên quan}
% \label{sec:related-work}
%
% \subsection{Constructive alignment và các khung đo alignment}
%
% Constructive alignment định nghĩa chất lượng giảng dạy qua quan hệ giữa mục tiêu
% học tập, hoạt động dạy--học và assessment, thay vì qua từng artifact riêng
% lẻ~\cite{biggs1996enhancing,biggs2011teaching}. Quality Matters đưa các quan hệ này
% vào quy trình review: objective phải đo được, còn assessment, học liệu và hoạt động
% phải hỗ trợ objective~\cite{qm2018rubric}. Về phía đo lường định lượng, chỉ số
% Porter đo mức tương hợp giữa hai phân bố nội dung--cognitive demand, các tiêu chí
% của Webb đánh giá những mặt như tính nhất quán của cognitive demand, và Fulmer khảo
% sát cách diễn giải các chỉ số
% này~\cite{porter2002measuring,webb1997criteria,fulmer2011estimating}. Bloom
% taxonomy bản sửa đổi cung cấp bộ từ vựng chung về cognitive
% demand~\cite{anderson2001taxonomy}. Điểm chung: các khung này đo trên curriculum,
% bộ chuẩn và assessment do con người soạn sẵn từ trước.
%
% \subsection{Hệ tác tử ISD và benchmark đánh giá}
%
% ISD-Agent-Bench tập hợp các agent xây theo những lý thuyết thiết kế khác nhau, từ
% ADDIE, thiết kế hệ thống, rapid prototyping tới vòng lặp
% reasoning--acting~\cite{isdagentbench,yao2023react}. Nhiều hệ khác hỗ trợ soạn
% giáo án theo định hướng sư phạm hoặc sinh lesson plan bằng cơ chế đa tác
% tử~\cite{lessonplanner2024,hu2024teachingplan,zhang2025eduplanner,agentlesson2026,zeng2026coma}.
% Về đánh giá, ISD-Agent-Bench và các benchmark giáo dục rộng hơn đều chấm bằng
% rubric LLM-judge~\cite{edubench2025,tutorbench2025,mathtutorbench2025}, trong khi
% bias của LLM-judge đã được ghi nhận có hệ
% thống~\cite{zheng2023judging,panickssery2024llm,norman2026reliability}. Xuyên suốt
% nhóm công trình này, chưa hệ nào đặt duy trì chất lượng khi thu nhỏ model làm mục
% tiêu trung tâm, và chưa benchmark nào instrument riêng các quan hệ giữa objective,
% hoạt động và assessment do máy sinh.
%
% \subsection{Các tiền lệ gần nhất}
%
% Về phía triển khai, các công trình gần nhất mỗi công trình giải quyết một phần bài
% toán. Small Models, Big Support chạy model self-hosted để tạo nội dung và
% assessment cho giáo viên, nhưng không có quy trình ISD agentic đầu
% cuối~\cite{reza2025smallmodels}. Jaldi và cộng sự so sánh SLM với LLM cho khâu sinh
% câu hỏi đánh giá với động cơ triển khai bảo vệ riêng tư, nhưng dừng ở một thành
% phần đơn lẻ~\cite{jaldi2026slmassessment}. LessonPlanLM đạt chất lượng lesson plan
% trên model 7B bằng fine-tune, thay vì kiểm tra một harness qua nhiều kích
% thước~\cite{zheng2025lessonplanlm}. Instructional Agents sinh trọn bộ học liệu và
% báo cáo chi phí, nhưng vẫn chạy trên API thương
% mại~\cite{yao2025instructionalagents}. EduClaw-Bench khảo sát tương tác giữa thiết
% kế agent và tier model, nhưng cho bài toán tutoring và vẫn chấm bằng
% LLM-judge~\cite{lee2026educlaw}. Như vậy, authoring tại chỗ, chuyên biệt hoá model
% nhỏ và sinh học liệu đa tác tử đã được nghiên cứu riêng lẻ~\cite{wei2026slmeducation};
% chúng chưa được kết hợp thành một phép đo có kiểm soát cho một ISD harness duy
% nhất trên dải model open-weight.
%
% Về phía đo lường, các thành phần cần thiết đều đã tồn tại. Semantic traceability và
% embedding đa ngôn ngữ khôi phục quan hệ giữa các artifact văn
% bản~\cite{clelandhuang2018traceability,feng2022labse}. Embedding và LM được dùng để
% phân loại assessment item vào bộ chuẩn nội dung có
% sẵn~\cite{xu2025itemalignment}. LessonBench-V1 so phân bố Bloom của bài giảng máy
% sinh với lesson plan tham chiếu đã qua duyệt của con
% người~\cite{silva2026lessonbench}. Với cognitive demand, Li và cộng sự công bố
% corpus learning objective do chuyên gia gán nhãn; nghiên cứu này dùng corpus đó để
% fine-tune Bloom classifier~\cite{li2022bloomclassification}, còn BloomNet là một
% classifier liên quan chứ không phải nguồn corpus~\cite{waheed2021bloomnet}. Các
% phương pháp trên không cần generative judge, nhưng đều đo sự ăn khớp với một bộ
% chuẩn hoặc tham chiếu có sẵn; chưa phương pháp nào đánh giá đồng thời các artifact
% do chính hệ bị đo sinh ra.
%
% \subsection{Tổng hợp}
%
% \Cref{tab:related} định vị các nhóm công trình theo tiêu chí; bảng mô tả phạm vi,
% không xếp hạng chất lượng. Hai tổ hợp còn bỏ ngỏ hiện rõ: đánh giá có kiểm soát
% theo scale cho một ISD harness chạy tại chỗ, và thước đo alignment nội tại tích
% hợp vào benchmark ISD --- nghiên cứu này xử lý tổ hợp thứ nhất bằng ladder Qwen
% 0.8B--9B và tổ hợp thứ hai bằng giao thức panel bảy tín hiệu.
%
% \begin{table}[t]
%   \caption{Định vị công trình liên quan theo tiêu chí. ISD = instructional design
%   agentic đầu cuối; Local = SLM local/open-weight là điều kiện thiết kế; Scale = đánh
%   giá có kiểm soát qua các kích thước model; Intr. = alignment nội tại giữa các
%   design artifact đều do máy sinh; Bench = tích hợp vào ISD benchmark.
%   \ding{51} = đã xét; \ding{55} = chưa xét; (\ding{51}) = đã xét một phần.}
%   \label{tab:related}
%   \centering
%   \small
%   \resizebox{\columnwidth}{!}{%
%   \begin{tabular}{@{}lccccc@{}}
%     \toprule
%     Công trình hoặc nhóm & ISD & Local & Scale & Intr. & Bench \\
%     \midrule
%     ISD-Agent-Bench và các agent~\cite{isdagentbench} & \ding{51} & \ding{55} & \ding{55} & \ding{55} & \ding{55} \\
%     Các hệ sinh lesson plan bằng LLM~\cite{lessonplanner2024,hu2024teachingplan,zhang2025eduplanner,agentlesson2026,zeng2026coma} & (\ding{51}) & \ding{55} & \ding{55} & \ding{55} & \ding{55} \\
%     Small Models, Big Support~\cite{reza2025smallmodels} & \ding{55} & \ding{51} & \ding{55} & \ding{55} & \ding{55} \\
%     SLM cho sinh câu hỏi đánh giá~\cite{jaldi2026slmassessment} & \ding{55} & \ding{51} & (\ding{51}) & \ding{55} & \ding{55} \\
%     LessonPlanLM~\cite{zheng2025lessonplanlm} & \ding{55} & (\ding{51}) & \ding{55} & \ding{55} & \ding{55} \\
%     Instructional Agents~\cite{yao2025instructionalagents} & (\ding{51}) & \ding{55} & \ding{55} & \ding{55} & \ding{55} \\
%     EduClaw-Bench~\cite{lee2026educlaw} & \ding{55} & \ding{55} & (\ding{51}) & \ding{55} & \ding{55} \\
%     Đo item--standards bằng LM/embedding~\cite{xu2025itemalignment} & \ding{55} & \ding{55} & \ding{55} & \ding{55} & \ding{55} \\
%     LessonBench-V1~\cite{silva2026lessonbench} & (\ding{51}) & \ding{55} & \ding{55} & \ding{55} & (\ding{51}) \\
%     Bloom classifier dựa trên BERT~\cite{li2022bloomclassification,waheed2021bloomnet} & \ding{55} & \ding{55} & \ding{55} & \ding{55} & \ding{55} \\
%     \textbf{AlignmentGraph-ISD} (nghiên cứu này) & \ding{51} & \ding{51} & \ding{51} & \ding{51} & \ding{51} \\
%     \bottomrule
%   \end{tabular}}
% \end{table}
